% integracion_frontend.tex — Capítulo 8: Integración con Frontend
\chapter{Integración con Frontend}

\section{Contexto de Integración}

El backend FINARA es consumido actualmente por la \textbf{app móvil Flutter} desarrollada en paralelo por el equipo TT 2026-A062. La plataforma web React está planificada para una fase futura. Ambos clientes consumen la misma API REST mediante solicitudes HTTP estándar con JSON.

\section{Configuración CORS}

El middleware CORS del backend permite solicitudes desde los orígenes configurados mediante la variable de entorno \texttt{ALLOWED\_ORIGINS}. Durante desarrollo, se acepta cualquier origen (\texttt{*}) para facilitar las pruebas locales.

\begin{lstlisting}[style=javascript, caption={Configuración CORS en app.js}]
app.use(cors({
  origin: process.env.ALLOWED_ORIGINS
    ? process.env.ALLOWED_ORIGINS.split(',')
    : '*',
  methods: ['GET', 'POST', 'PUT', 'DELETE', 'OPTIONS'],
  allowedHeaders: ['Content-Type', 'Authorization'],
}));
\end{lstlisting}

\begin{lstlisting}[style=bash, caption={Variable de entorno para producción}]
# .env (produccion)
ALLOWED_ORIGINS=https://finara.app,app://finara
\end{lstlisting}

\section{Consumo desde la App Móvil Flutter}

La app Flutter implementa un servicio \texttt{ApiService} que encapsula todas las llamadas al backend. Los ejemplos siguientes muestran los patrones de integración más representativos.

\subsection{Autenticación y Almacenamiento del Token}

\begin{lstlisting}[style=bash, caption={Login y almacenamiento del JWT en Flutter (Dart)}]
// En Flutter (Dart) - auth_service.dart
Future<LoginResponse> login(String email, String password) async {
  final response = await http.post(
    Uri.parse('http://192.168.1.x:3000/api/auth/login'),
    headers: {'Content-Type': 'application/json'},
    body: jsonEncode({'email': email, 'password': password}),
  );

  if (response.statusCode == 200) {
    final data = jsonDecode(response.body);
    final token = data['data']['token'];

    // Almacenar token de forma segura
    await FlutterSecureStorage().write(key: 'jwt_token', value: token);

    return LoginResponse.fromJson(data['data']);
  }
  throw ApiException(jsonDecode(response.body)['error']);
}
\end{lstlisting}

\subsection{Request Autenticado Estándar}

\begin{lstlisting}[style=bash, caption={Patrón de request autenticado en Flutter}]
// Obtener gastos del mes actual
Future<List<Expense>> getExpenses({int? month, int? year}) async {
  final token = await FlutterSecureStorage().read(key: 'jwt_token');
  final now   = DateTime.now();

  final uri = Uri.parse(
    'http://192.168.1.x:3000/api/expenses'
    '?month=${month ?? now.month}&year=${year ?? now.year}'
  );

  final response = await http.get(uri, headers: {
    'Authorization': 'Bearer $token',
    'Content-Type':  'application/json',
  });

  if (response.statusCode == 200) {
    final data = jsonDecode(response.body)['data'] as List;
    return data.map((e) => Expense.fromJson(e)).toList();
  }
  throw ApiException(jsonDecode(response.body)['error']);
}
\end{lstlisting}

\subsection{Subida de Imagen para OCR}

\begin{lstlisting}[style=bash, caption={Envío de imagen de ticket para escaneo OCR}]
// ticket_service.dart
Future<TicketScanResult> scanTicket(File imageFile) async {
  final token   = await FlutterSecureStorage().read(key: 'jwt_token');
  final request = http.MultipartRequest(
    'POST',
    Uri.parse('http://192.168.1.x:3000/api/tickets/scan'),
  );

  request.headers['Authorization'] = 'Bearer $token';
  request.files.add(
    await http.MultipartFile.fromPath('images', imageFile.path)
  );

  final streamedResponse = await request.send();
  final response         = await http.Response.fromStream(streamedResponse);

  if (response.statusCode == 201) {
    return TicketScanResult.fromJson(
      jsonDecode(response.body)['data']
    );
  }
  throw ApiException(jsonDecode(response.body)['error']);
}
\end{lstlisting}

\section{Manejo de Errores desde el Cliente}

La Tabla \ref{tab:http-errors} describe cómo la app Flutter interpreta cada código de respuesta del backend.

\begin{longtable}{>{\centering}p{1.5cm} p{4cm} p{7cm}}
  \caption{Códigos HTTP de la API y su interpretación en Flutter}
  \label{tab:http-errors} \\
  \toprule
  \textbf{Código} & \textbf{Significado} & \textbf{Acción recomendada en Flutter} \\
  \midrule
  \endfirsthead
  \toprule
  \textbf{Código} & \textbf{Significado} & \textbf{Acción recomendada en Flutter} \\
  \midrule
  \endhead
  200 & Éxito                    & Procesar \texttt{data} del JSON \\
  201 & Recurso creado           & Actualizar lista, navegar a detalle \\
  400 & Datos inválidos          & Mostrar \texttt{error} del JSON en el formulario \\
  401 & No autenticado / Token expirado & Borrar token, redirigir a Login \\
  404 & Recurso no encontrado    & Mostrar mensaje ``No encontrado'' \\
  429 & Demasiadas solicitudes   & Esperar y reintentar (leer \texttt{Retry-After}) \\
  500 & Error del servidor       & Mostrar mensaje genérico de error \\
  \bottomrule
\end{longtable}

\section{Formato de Fechas}

El backend retorna fechas en formato ISO 8601 (\texttt{YYYY-MM-DDTHH:mm:ss.sssZ}) y acepta fechas en formato \texttt{YYYY-MM-DD} en los request bodies. La app Flutter utiliza el paquete \texttt{intl} para formatear estas fechas según la localidad del usuario (es\_MX).

\section{URL Base de la API}

\begin{longtable}{p{3cm} p{10.5cm}}
  \caption{URL base según entorno}
  \label{tab:base-url} \\
  \toprule
  \textbf{Entorno} & \textbf{URL Base} \\
  \midrule
  Desarrollo local  & \texttt{http://10.0.2.2:3000/api} (emulador Android $\to$ localhost) \\
  Dispositivo físico & \texttt{http://192.168.x.x:3000/api} (IP local de la PC en la red WiFi) \\
  Producción        & \texttt{https://api.finara.app/api} (dominio de producción) \\
  \bottomrule
\end{longtable}

\begin{notaTecnica}
  En el emulador Android, \texttt{localhost} del dispositivo es diferente al \texttt{localhost} de la PC anfitriona. Para acceder al backend corriendo en la PC, el emulador usa la IP especial \texttt{10.0.2.2}.
\end{notaTecnica}
